\section{Background}
\subsection{Hardware-side-channel defenses}\label{sec:bg-defenses}
Many HSCDs exist that include some hardware component, as viable performance necessitates HSCAs to be mitigated at least partially in hardware. These defenses aim to protect some set of architectural state from being leaked through some side channel formed by a data-dependent optimization. Many of these defenses function by designating a set of architectural state that contains secrets, and tracking how information flows from an secret piece of state to a speculative side channel. STT, or Speculative Taint Tracking, is a defense that stalls transmitters, or instructions that create data-dependent behavior, until they are non-speculative when they depend on a tainted, or secret, piece of state \cite{yu2019stt}.
\notebr
\auditnote{cite}{proposed}{\textbf{ADD (new sentence, SPT; insert here, before ReCon)}: ``SPT~\cite{spt} delays every transmitter until it can prove that the transmitter's operands are leaked by the program's non-speculative execution anyway, using a dynamic information-flow analysis to declassify data as it leaks.'' \newline \textbf{REWORD (optional, the STT sentence above)}: end it with ``\ldots until the instruction that produced the tainted data becomes non-speculative, i.e., reaches its visibility point~\cite{yu2019stt}'', which sets up the visibility-point race in the results. \newline \textit{Source:} choudhary-2021-spt.pdf abstract (``delays execution of every transmitter until it can prove that the transmitter's operands leak during the program's non-speculative execution''); yu-2019-stt.pdf abstract (``the earliest time is when the instruction(s) producing the protected data become non-speculative'').}
\notebr
ReCon builds on such defenses by shrinking what they protect: it detects data that has already leaked during non-speculative execution, limiting detection to pairs of directly dependent loads that dereference a pointer, records that knowledge in the coherence protocol, and lets a defense such as STT treat that data as safe to use under speculation \cite{recon}.
\notebr
\reviewnote{Anna, check this ReCon sentence: it replaces the NDA sentence (as you asked) and is paraphrased from ReCon's abstract in references.bib.}
\notebr
Many more hardware HSCDs exist, some of which improve the performance of STT by actively eliminating items from the STT protection set. All of these defenses protect against different combinations of leakage components, including different protection sets, and different types of side channels.
\notebr
\auditnote{cite}{proposed}{\textbf{ADD (new sentences; insert after this paragraph, or in place of its last sentence)}: ``Other defenses let speculative loads execute and instead control their effect on the cache: InvisiSpec~\cite{yan2018invisispec} keeps speculative loads invisible to the cache hierarchy and exposes each load only once it is safe, CleanupSpec~\cite{saileshwar2019cleanupspec} lets speculative loads change the cache and undoes those changes on a mis-speculation, and SpecLFB~\cite{cheng2024speclfb} holds speculative misses in the cache's line-fill buffer and installs them only once they are safe. All of these defenses target Spectre attacks~\cite{spectre}, including store bypass (Spectre v4), in which a load speculatively reads stale data before an older store's address is known~\cite{horn2018ssb}.'' \newline \textit{Source:} fu-2025-amulet.pdf \S4.5 (InvisiSpec: ``speculative loads are invisible to the caches; when loads become safe, an expose operation makes the load visible''), \S4.6 (CleanupSpec: ``allows speculative loads to modify cache state and cleans up these state changes on a mis-speculation''), \S4.7 (SpecLFB: ``adds security checks to the cache Line-Fill Buffer\ldots blocks speculative cache misses from being installed into the cache until they are safe''); the store-bypass wording follows your own pipeline.tex description (check it against horn2018ssb). \newline \textbf{Two caveats (2026-11-09):} keep the SpecLFB clause as a description of the paper's design only: the gem5 SpecLFB build bounces unsafe misses and replays them, and its park-in-the-LFB path is unreachable, so do not carry it into the Results. Our STT tests are branch-only and the store-bypass tests run only on the Kaxiras-group builds, so write ``All of these designs target \ldots'' rather than implying that every defense is tested for store bypass. \textit{Source:} simspect URGENT\_TODO.md (``SpecLFB parks unsafe misses in the LFB'' correction); wiki/state/EVAL\_MANIFEST.md STL section.}
\urgentbr
\urgentnote{\textbf{VOCAB}: \textbf{ADD (new paragraph: leakage contracts)} here: this is where you separate our terms from Guarnieri et al.'s and AMuLeT's. \textbf{PROPOSED TEXT}: ``Guarnieri et al.\ specify what a CPU may leak as a \emph{hardware--software contract}~\cite{guarnieri2021contracts}: an \emph{observation clause} says what each executed instruction exposes to an attacker (for example, the addresses of loads and stores), and an \emph{execution clause} says which instructions may execute speculatively. Revizor and AMuLeT test CPUs and simulators against such contracts~\cite{revizor,amulet}, which Revizor calls speculation contracts and AMuLeT calls leakage contracts. We use \emph{leakage contract} in the narrower sense of SynthLC~\cite{rtl2upath}: the set of (transmitter, operand) pairs through which a design can leak. Our leakage contract is therefore closest to their observation clause, though it names the operands that are fundamentally leaked rather than the attacker-visible state; our execution contract is closest to their execution clause; and our protection set makes explicit which data must not reach those operands, which a contract leaves implicit (everything its traces do not expose).'' Check the last clause: it is an interpretation of contract satisfaction, not a statement from the papers. Also cite here: Spectector~\cite{guarnieri2020spectector} for SNI, LCM~\cite{lcms}, and the RISC-V follow-ups~\cite{mohr2024synthesizing,synthlc,leave}. \textit{Source:} fu-2025-amulet.pdf Sec.~2.3 (``A leakage contract captures the expected leakage of a CPU under test'') and Table~1 (CT-SEQ's observation and execution clauses); Revizor's title; RTL2M$\mu$PATH (MICRO 2024) on SynthLC's ``formally verified leakage contracts''.}
\notebr

\subsection{HSCD litmus testing}\label{sec:bg-litmus}
Prior work that validates advertised security guarantees against hardware support for hardware SCDs analyzes fabricated designs\cite{revizor} or RTL, due to the wide availability\cite{rtl2upath}.
\notebr
\auditnote{cite}{cite, high}{The sentence before this note still ends in the unfinished clause `due to the wide availability', and Revizor tests commercial black-box CPUs against speculation contracts, not a defense's advertised guarantee. The intro's prior-work sentence (``Existing validation either tests fabricated CPUs\ldots'') now says this with the full cite list, so point back to it or reuse it here. \newline \textbf{PROPOSED TEXT} (edits the sentence just before this note): Prior work validates side-channel security on fabricated CPUs~\cite{revizor,oleksenko2023hideseekspectresefficient} and on RTL designs~\cite{rtl2upath,contractshadow,hur2022specdoctor}. \newline \textit{Source:} references.bib entries revizor, oleksenko2023hideseekspectresefficient (S\&P 2023), rtl2upath (MICRO 2024), contractshadow (ASPLOS 2025), hur2022specdoctor (CCS 2022); sections/intro.tex prior-work sentence.}
\notebr
Only AMuLeT\cite{amulet} analyzes simulator implementations. AMuLeT is a fuzzer that generates validation tests to detect security bugs in simulated Spectre defense prototypes—one of our goals. It does random test generation and conducts relational testing by running programs in simulation that, when supplied with two different inputs, are indistinguishable in their architectural side-channel leakage, flagging differentiated microarchitectural state. However, it attempts to validate high-level security properties, and has no ability to do directed testing for leakage.

\anna{talk abt uarch state}

\subsection{Axiomatic execution models}\label{sec:bg-axiomatic}
Prior work has formalized processor memory consistency models (MCMs) into both axiomatic and operational models \cite{herdingcats}\anna{cite a bunch of other stuff, arm + simner stuff}. Axiomatic MCM models represent an execution as a graph of memory events connected by relations such as program order (po), which relates an instruction to its subsequent instruction in thread-order, reads-from (rf) which relates a write to a read that uses its value, and coherence (co), which relates two writes to the same address. Lustig et. al. use this format of model along with a set of minimality requirements to generate a comprehensive set of MCM-violating execution graphs and concretize them into litmus tests for testing an MCM implementation \cite{lustig}. They do this by using the Alloy Analyzer \cite{jackson2002alloy} to generate all tests that fail according to their MCM and are minimal
\notebr
\auditnote{cite}{cite, high}{The sentence before this note has no final period, and Fig.~3 names the Kodkod model finder, which has no bib entry yet (jackson2002alloy is now cited here). \newline \textbf{PROPOSED TEXT} (edits the sentence just before this note): They use the Alloy Analyzer~\cite{jackson2002alloy,kodkod} to generate every minimal test that their MCM forbids. New bib entry \texttt{kodkod}: E.~Torlak and D.~Jackson, ``Kodkod: A relational model finder,'' TACAS 2007 (as cited in lustig-2017-automated-litmus-synthesis.pdf and trippel-2018-checkmate.pdf ref.~[33]).}
\notebr
\anna{replace this sentence, cite wickerson}. Minimality, as defined in the test generation procedure, means that a graph violates the desired predicate, and cannot be relaxed, i.e. have a component removed or demoted, without becoming a non-violating graph. The relaxations considered by Lustig et. al. include removing instructions and dependencies, as well as demoting fences. Minimality allows prioritization of axes for test diversity.

Axiomatic modeling has been used for formalization of speculation restrictions \anna{cite checkmate, all of the speculation formalization, LCM}. 

\anna{add way more about modeling executions}
% A similar such axiomatic definition of side channel leakage would allow generation of a similar test suite satisfying a minimality criterion. 
\notebr
\auditnote{cite}{proposed}{\textbf{WRITE MORE HERE (axiomatic modeling)}: write this out more here (you said you would). Proposed thread and citations: memory models as relations over executions~\cite{herdingcats,alglave2010fences,batty2011mathematizing,sarkar2011power,pulte2018simplifying,lahav2017repairing}; Alloy and relational model finding~\cite{jackson2002alloy}; comparing and synthesizing models and litmus tests~\cite{memalloy,lustig,memsynth,alglave2011litmus,mcmutants}; microarchitectural happens-before ($\mu$hb) models~\cite{lustig2014pipecheck,manerkar2015ccicheck,lustig2016coatcheck,trippel2017tricheck,lustig2015armor,hossain2020transform}; full ISA and system semantics~\cite{simner2020ifetch,simner2022vm,simner2025exceptions,armstrong2021isla,armstrong2019sail}; and the step to security: CheckMate's $\mu$hb exploit patterns~\cite{checkmate}, Cats vs.\ Spectre~\cite{ponce2022cats} and LCM~\cite{lcms}. End with what SimSpect takes from this line of work (execution graphs with rf/spo-style relations, minimal litmus enumeration) and what it adds (speculation flags and a defense's leakage predicate). Since 2026-11-08 the enumeration has two stages, so say ``minimal litmus enumeration (Stage~1a, RI-minimal graphs), then a capped flesh-out under loosened minimality (Stage~1b)'' rather than only ``minimal litmus enumeration''.}
\notebr

